% 제출본 전용 — 3장 끝에 붙는다 (3.8). 공학인증 설계 구성요소 표.
% main_ko.tex 는 이 파일을 읽지 않는다.

\subsection{설계 구성요소 연관성 --- 설계 및 제작}

이 작품의 산출물은 화면이 있는 응용 프로그램이 아니라 \textbf{분석
파이프라인과 그 위에서 내린 설계 판단}이다. 공학 설계의 구성요소가 이 작품에서
무엇에 해당하는지를 밝힌다.

\begin{table}[htbp]
  \centering
  \caption{설계 구성요소 연관성 --- 설계 및 제작}
  \label{tab:v2-design}
  \begin{tabular}{llcp{9.0cm}}
    \toprule
    \textbf{설계 구성요소} & \textbf{구분} & \textbf{연관성} & \textbf{내용} \\
    \midrule
    \multirow{3}{*}{설계 구성요소}
      & 범위결정 & O &
        주 결과변수를 가격에서 \textbf{거래량}으로, 인과 추정 구간을
        2025-10-19 이전으로, 처치군을 동 전체가 지정된 4개 동으로 확정
        (0001 · 0002 · 0012). 무엇을 묻지 \emph{않을지}도 함께 정했다 \\
    \cmidrule(l){2-4}
      & 요구분석 & O &
        답해야 할 \textbf{다섯 질문}과 각 질문이 성립하기 위한 조건을
        \texttt{docs/framing.md}에 명세했다. 메커니즘에서 세운
        \textbf{사전 등록 예측}(4.7절)이 요구사항의 검증 기준이 된다 \\
    \cmidrule(l){2-4}
      & 설계 & O &
        설계 선택마다 근거와 \textbf{버린 대안}을 결정기록
        \pipeDecisions 편으로 남겼다. 바뀐 판단은 지우지 않고 그 자리에
        정정으로 적어 변경 이력이 보존된다 (부록 A) \\
    \bottomrule
  \end{tabular}

  \begin{minipage}{\linewidth}\footnotesize
  \textit{주.} 괄호 안의 네 자리 수는 결정기록 번호다. 전문은 저장소의
  \texttt{docs/decisions/}에 있다\citep{repo}.
  \end{minipage}
\end{table}

\textbf{요구사항명세서와 설계문서가 따로 있는 것이 아니라 그 역할을 하는 문서가
있다.} 일반적인 소프트웨어 개발에서 요구사항명세서가 하는 일 --- 무엇을 만들
것인지, 무엇이 충족되면 된 것인지를 미리 적어 두는 일 --- 을 이 작품에서는
\texttt{framing.md}의 다섯 질문과 사전 등록 예측이 한다. 설계문서가 하는 일
--- 왜 그렇게 만들었는지를 남기는 일 --- 은 결정기록이 한다.

차이가 하나 있다. \textbf{설계문서는 보통 완성된 설계만 담지만, 결정기록은 버린
대안과 바뀐 판단까지 담는다.} 6장에서 이질성의 원인에 대한 서술이 세 번 바뀌는
과정을 그대로 실을 수 있는 것은 이 기록 때문이다.
